\documentclass[11pt]{article}
\usepackage[margin=1in]{geometry}
\usepackage{amsmath,amssymb,amsthm,mathtools}
\usepackage{microtype,hyperref}
\newtheorem{theorem}{Theorem}
\newtheorem{proposition}[theorem]{Proposition}
\newtheorem{lemma}[theorem]{Lemma}
\newtheorem{corollary}[theorem]{Corollary}
\theoremstyle{remark}\newtheorem{remark}[theorem]{Remark}
\newcommand{\ddc}{dd^c}
\newcommand{\dA}{\mathrm dA}
\newcommand{\PhiG}{\Phi}
\newcommand{\CP}{\mathbb P}
\newcommand{\D}{\mathbb D}
\newcommand{\R}{\mathbb R}
\newcommand{\C}{\mathbb C}
\DeclareMathOperator{\Rea}{Re}
\title{Growing vanishing jets on compact curves:\\
the universal range and a mesoscopic free-boundary correction}
\author{}
\date{3 October 2026}
\begin{document}
\maketitle
\begin{abstract}
For partial Bergman kernels selected by a growing vanishing order on a
compact curve with local weight $|z|^{2r}+O(|z|^{2r+1})$, we prove
the optimal universal range of the uncorrected Gaussian transition as
$q=o(m^{1/(r+1)})$. A correction for every $q\to\infty$,
$q=o(m)$ centers the transition on a small curvature quadrature domain
and uses the normal curvature scale. The argument includes a direct
near-radial free-boundary construction, Gaussian annulus kernels and a
singular-metric correction preserving the growing ideal. First critical
Taylor corrections and compact counterexamples establish the optimality
of the range and the limitations of finite Taylor information.
\end{abstract}

\section{Data and statements}
Let $C$ be a connected compact Riemann surface, $x\in C$, and $(L,h)$ a
smooth Hermitian holomorphic line bundle. Assume that its curvature is
semipositive and is strictly positive on $C\setminus\{x\}$. In a
coordinate $z$ centered at $x$ and a holomorphic frame $e$, suppose
\[
 |e|_h^2=e^{-\phi},\qquad
 \phi(z)=|z|^{2r}+O(|z|^{2r+1}),\qquad r\in\mathbb Z_{\geq1}.
\]
Fix a smooth positive volume form $\nu=a(z)\dA$, with $a(0)=1$ locally,
where $\dA$ is Euclidean area. Our convention is
\[
 \ddc=\frac{i}{2\pi}\partial\bar\partial,
 \qquad \rho=\phi_{z\bar z},\qquad c_1(L,h)=\frac{\rho}{\pi}\dA.
\]
Let $q=q_m$ be integers such that $q\to\infty$ and $q=o(m)$. Give
\[
 V_{m,q}=H^0(C,L^m\otimes\mathfrak m_x^q)
 \quad\hbox{the norm}\quad
 \|s\|^2=\int_C|s|_{h^m}^2\nu.
\]
Write $B_{m,q}=\sum_j|s_j|_{h^m}^2$ for an orthonormal basis, and set it
to zero if the space is zero. Put
\[
 t=q/m,\qquad R=(t/r)^{1/(2r)},\qquad
 \PhiG(u)=\frac1{\sqrt{2\pi}}\int_{-\infty}^u e^{-v^2/2}\,\mathrm dv.
\]
Every asymptotic statement below is for fixed geometric data and fixed
$r$; constants may depend on them.
These are the data of catalogue P22~\cite{catalogue}. The known compact
small-growth theorem is Theorem 26.13.1 of Xia's working
monograph~\cite{book}; its statement is used only as background.
The degree is the integral of the Chern curvature, so $\deg L>0$.
Riemann--Roch~\cite[Chapter VI, Section 10]{DA} gives
$\dim V_{m,q}\geq m\deg L-q+1-g>0$ eventually, where $g$ is the
genus. Thus the zero-space convention has no asymptotic effect.

\begin{theorem}[Corrected transition]\label{thm:corrected}
For all small $t>0$ there is a simply connected domain $D_t$ near $x$,
with boundary $\Gamma_t$, such that $R^{-1}D_t$ tends to the unit disk
and
\begin{equation}\label{eq:quadrature}
 \int_{D_t}\rho(z)f(z)\dA=\pi t f(0)
\end{equation}
for every $f$ holomorphic near its closure. Let $d_t(z)$ be signed
Euclidean distance from $\Gamma_t$, positive on the exterior, and let
$p_t(z)$ be the nearest boundary point. For every fixed $M$, uniformly
in all directions and $|d_t(z)|\leq MR/\sqrt q$,
\begin{equation}\label{eq:edge}
 \frac{\pi a(z)B_{m,q}(z)}{m\rho(z)}
 -\PhiG\bigl(2\sqrt{m\rho(p_t(z))}\,d_t(z)\bigr)\longrightarrow0.
\end{equation}
If $d_t(z)=o(R)$ and $\sqrt q\,d_t(z)/R$ tends to $-\infty$ or
$+\infty$, the same normalized kernel tends to $0$ or $1$, respectively.
Thus the corrected center and transition length are $\Gamma_t$ and
$1/(2\sqrt{m\rho})$, for the entire sublinear range.
\end{theorem}

\begin{theorem}[Universal range]\label{thm:range}
For every $r\geq1$, the law
\begin{equation}\label{eq:raw}
 \frac{\pi B_{m,q}(z_m)}{mr^2|z_m|^{2r-2}}
 \longrightarrow\PhiG(\tau),\qquad
 m|z_m|^{2r}=q/r+\tau\sqrt{q/r},
\end{equation}
holds for every permissible metric and every sequence of observation
directions when
$q=o(m^{1/(r+1)})$. This is the optimal universal range: for every
excluded sequence $q\to\infty$, $q=o(m)$, a compact example satisfying
the same hypotheses violates the law along a subsequence.
\end{theorem}

Write $\Gamma_t$ as $|z|=R c_R(\theta)$. For a fixed metric, the exact
criterion for \eqref{eq:raw} for every sequence of directions and every
fixed $\tau$ is
\begin{equation}\label{eq:criterion}
 \sqrt q\,\|c_R-1\|_\infty\longrightarrow0.
\end{equation}
The first critical correction and further effective coefficients are
derived in Section~\ref{sec:corrections}.

\section{A direct near-radial free-boundary construction}
The geometric input is proved here for the particular rescaled family.
General spatial regularity of obstacle boundaries alone does not imply
arbitrary smooth parameter dependence. For example,
Theorem~1.1 of Serfaty--Serra \cite{SS} establishes two parameter
derivatives and discusses stronger joint regularity as an open question.
The following argument instead uses a fixed-boundary Fourier equation.

\begin{lemma}\label{lem:hole}
Let $Q(R,\zeta)$ be jointly smooth on a fixed disk of radius $b>1$,
with $Q(0,\zeta)=|\zeta|^{2r}$. Suppose every $Q_R$ is subharmonic
and $Q_{R,\zeta\bar\zeta}>0$ in a fixed collar of the unit circle.
For small $R$, the largest subharmonic minorant of $Q_R$ with a pole
$2r\log|\zeta|+O(1)$ at zero has exactly one noncontact component
$\Omega_R$. Its boundary is a radial graph near the unit circle. For
every fixed finite $k$ the normalized interior map $F_R$ and the
logarithm-subtracted harmonic part depend smoothly on $R$ in the
corresponding H\"older spaces, with $O(R)$ distance from their radial
values. All estimates used below require only finitely many such norms.
\end{lemma}
\begin{proof}
Fix $k\geq3$ and $0<\alpha<1$. Let $A^{k,\alpha}$ be the real Banach
space of functions holomorphic on $\D$, $C^{k,\alpha}$ on its closure,
whose value at zero is real. Set $f(u)=u(1+h(u))$. Small $h$ gives a
univalent $f$, with $f(0)=0$ and $f'(0)>0$. On $|u|=1$ define
\[
 K(R,h)=f(u)(Q_R)_\zeta(f(u))-r,
 \qquad T(R,h)=(P_-K,\Rea[K]_0),
\]
where $P_-$ projects onto strictly negative Fourier modes and $[\cdot]_0$
is the constant coefficient. The boundedness of Cauchy projections on
these H\"older spaces follows from the circle Hilbert kernel: subtract
the value at the evaluation point, split its integral at twice the
increment length, and use the H\"older difference bound on the inner
part and the differentiated kernel bound on the outer part. This gives
the $C^{0,\alpha}$ bound; commuting angular derivatives gives the
$C^{k,\alpha}$ bound; see also \cite[Theorem 1.2]{WD}.
The smooth superposition operator has no derivatives
of $h$. Thus $T$ is a smooth map to the negative-mode space times $\R$.
At $(R,h)=(0,0)$, $T=0$, and
\[
 D_hK(0,0)v=r^2(v+\bar v),\qquad
 D_hT(0,0)v=(r^2\overline{v-v(0)},\,2r^2v(0)).
\]
This is a bounded real-linear isomorphism. The Banach implicit-function
theorem~\cite[Theorem 2]{K} therefore gives a unique normalized branch $f_R$, smoothly
depending on $R$ in the fixed space, with $\|f_R-\mathrm{id}\|=O(R)$.
Branches obtained in higher fixed spaces agree by lower-space uniqueness.
This supplies all finite spatial estimates and all finite parameter
expansions needed in the asymptotic argument. Each prescribed finite
order has its own sufficiently small parameter interval; no common
interval is inferred by intersecting infinitely many intervals.
Boundary regularity can
also be viewed as a scalar Riemann--Hilbert problem on the totally real
graph $\{(\zeta,\zeta(Q_R)_\zeta-r)\}$; its antiholomorphic derivative
is $\zeta(Q_R)_{\zeta\bar\zeta}\ne0$ near the boundary. Related smooth
disk-deformation frameworks are given in \cite{F,L,RWN}.

The equation has not omitted an imaginary residue condition. The
function $K$ now extends holomorphically to $\D$, and $K(0)$ is purely
imaginary. Stokes' formula gives
\[
 \frac1{2\pi i}\int_{|u|=1}(Q_R)_\zeta(f_R(u))f_R'(u)\,\mathrm du
 =\frac1\pi\int_{f_R(\D)}(Q_R)_{\zeta\bar\zeta}\dA\in\R.
\]
The residue of $(r+K)f_R'/f_R$ is $r+K(0)$, so $K(0)$ is real as well,
and is zero. Consequently
\[
 S_R(\zeta)=\frac{r+K(f_R^{-1}(\zeta))}{\zeta}
 =r/\zeta+H_R(\zeta)
\]
has a holomorphic $H_R$ on $\Omega_R=f_R(\D)$ and equals $(Q_R)_\zeta$
on its boundary. Let
\[
 U_R=2r\log|\zeta|+2\Rea\int_0^\zeta H_R(w)\,\mathrm dw+C_R.
\]
Its complex gradient agrees with that of $Q_R$ on the boundary.
The boundary difference is constant; choose $C_R$ to make it zero.
The difference and its gradient vanish there. Its inward normal second
derivative is $\Delta Q_R>0$, because the tangential second derivative
vanishes. Hence $Q_R-U_R>0$ in a uniform inner collar. At $R=0$,
\[
 U_0=1+2r\log|\zeta|,\qquad
 |\zeta|^{2r}-1-2r\log|\zeta|>0\quad(0<|\zeta|<1).
\]
Continuity preserves this inequality on the remaining compact interior
sets; the logarithmic term controls a neighborhood of zero.

Patch $U_R=Q_R$ off $\Omega_R$. Value and gradient matching eliminate
boundary measures in its Laplacian. Its curvature inside is $r\delta_0$
and outside is $\ddc Q_R\geq0$, so it is an admissible minorant. If $v$
is any competitor, $v-U_R$ is subharmonic on $\Omega_R\setminus\{0\}$,
bounded above at zero and nonpositive on its boundary. Removable
subharmonic extension and the maximum principle give $v\leq U_R$.
Outside the hole this follows from $v\leq Q_R$. This identifies the
envelope and excludes extra noncontact components. The radial graph and
the harmonic-part estimates follow from the normalized map and its
parameter dependence. We henceforth denote this map by $F_R=f_R$.
\end{proof}

Taylor's integral remainder gives, for small signed $R$,
\[
 Q_R(\zeta)=R^{-2r}\phi(R\zeta)
 =|\zeta|^{2r}+R E(R,\zeta)
\]
with jointly smooth $E$. This applies the lemma to the actual weight.
Write its envelope inside the hole as
\[
 U_R=2r\log|\zeta|+2\Rea G_R(\zeta),\qquad
 G_R=1/2+O(R).
\]
Set $D_t=R\Omega_R$. In physical coordinates the obstacle potential is
$\phi_t(z)=R^{2r}U_R(z/R)$ inside $D_t$ and $\phi$ outside.
It is a global singular metric on $L$, is $C^1$ across the boundary,
is bounded above by $\phi$, and has pole $2t\log|z|+O(1)$.
On the boundary $\phi_z=t/z+g_t'$ for a holomorphic $g_t$, so Green's
formula gives \eqref{eq:quadrature}. Its curvature is $t\delta_x$
inside and the original curvature outside.

\section{A correction estimate at the shrinking radius}
\begin{lemma}\label{lem:dbar}
There is a bounded smooth bump $b_R$ such that the singular potential
$\psi_m=m\phi_t+b_R$, with the inverse canonical metric included,
has curvature at least $c/R^2$ globally. On exterior annuli of radii
comparable to $R$ the lower bound is $cq/R^2$. If an $L^m$-valued
$\bar\partial$ source $\alpha$ is supported on such an exterior annulus,
it has a linear solution $v$ with
\begin{equation}\label{eq:dbar}
 \|v\|_{h^m}^2\leq C\frac{R^2}{q}\|\alpha\|_{h^m,\nu}^2.
\end{equation}
If the source vanishes near $x$, $v$ is holomorphic there and vanishes
to order at least $q$.
\end{lemma}
\begin{proof}
Take $b_R=|z|^2/R^2$ on a disk containing the hole, and cut it off before
radius $3R$, with its negative curvature supported outside the hole.
The bump is uniformly bounded and its Hessian is $O(R^{-2})$. On that
support the curvature $m\rho$ is comparable to $q/R^2$, so $q\to\infty$
absorbs the negative part. Inside the hole the bump supplies $c/R^2$.
Elsewhere in a fixed coordinate disk the original curvature is bounded
below by a fixed multiple of $m|z|^{2r-2}$; off that disk it is at least
$cm$. Since $mR^2\to\infty$, all these bounds dominate $c/R^2$ and the
bounded Ricci term of $\nu$.

Apply the pointwise inverse-curvature $L^2$ theorem for singular metrics,
as stated in Demailly \cite[Theorem~(5.2) and its following paragraph]{D}.
For sections use $L^m\otimes K_C^{-1}$ in bidegrees $(1,0)$ and $(1,1)$;
this includes the Ricci term and preserves the original form norms under
the canonical identification. The inverse curvature on the support is
$O(R^2/q)$. There $\phi_t=\phi$, and the bump is bounded. Elsewhere
$e^{-m\phi}\leq C e^{-\psi_m}$. These comparisons yield \eqref{eq:dbar}.
To use a continuous positive curvature lower bound, interpolate from
$c/R^2$ inside the hole to $cq/R^2$ across a thin exterior collar.
For the separated cutoff supports below, choose this collar to miss
the source. More generally, for an exterior source accumulating on
$\Gamma_t$, shrink the collar to the boundary. Dominated convergence
in the inverse-curvature integrals and weak compactness in the fixed
weighted Hilbert space give the same estimate; the boundary has area
zero. No discontinuous lower bound is required by the cited theorem.
Approximation of the positive singular metric gives the same estimate
and a weak solution. Distributional $\bar\partial$ regularity
\cite[Corollary 13.12]{W} makes it holomorphic where the source vanishes.
Near $x$, weighted integrability is against
$|z|^{-2q}$ times a smooth positive factor. A holomorphic term $z^j$ is
integrable precisely when $j\geq q$; thus the solution has the required
vanishing order.
\end{proof}

\section{Straightening and Gaussian annulus kernels}
Define on the unit disk
\[
 A_R(w)=G_R(F_R(w))+r\log(F_R(w)/w),\qquad
 \Psi_R=Q_R(F_R(w))-2\Rea A_R(w)-2r\log|w|.
\]
The logarithm uses the branch near one. Here $A_R=1/2+O(R)$ and
$F_R=w+O(R)$ in every fixed norm used below. The effective potential is
positive inside, with zero value and gradient on the boundary. For
$w=e^{v+i\theta}$ its expansion is
\begin{equation}\label{eq:ring}
 \Psi_R=(2r^2+O(R))v^2+O(v^3).
\end{equation}
The exact quadratic coefficient is
$2(Q_R)_{\zeta\bar\zeta}(F_R(e^{i\theta}))|F_R'(e^{i\theta})|^2$.

The maps need not extend holomorphically across the circle. To obtain
genuine exterior charts, truncate their Taylor series at
$K=\lfloor q^{1/16}\rfloor$, denoting them by $F_K,A_K$.
Smooth Fourier coefficient bounds in a sufficiently high fixed norm
make the truncation error $O(q^{-10})$, including the fixed derivatives
needed here. On $|w|\leq1+c/K$ their derivatives remain bounded and
$F_K-w,A_K-1/2=O(R)$: the factor $(1+c/K)^K\leq e^c$ controls the sums.
In particular $F_K$ is univalent there for small $R$. With
$\delta=c/K$ and
\[
 \Psi_K=Q_R(F_K)-2\Rea A_K-2r\log|w|,
\]
value and first derivative errors on the circle are $O(q^{-10})$.
Uniformly for $|v|\leq\delta$,
\begin{equation}\label{eq:ringK}
 \Psi_K=(2r^2+O(R)+O(|v|))v^2+O(q^{-10})(1+|v|).
\end{equation}
On the interior the effective potential has uniform positive quadratic
growth near the circle, a positive gap on compact annuli farther inside,
and logarithmic growth at zero.

In $z=RF_K(w)$ a local holomorphic section of order at least $q$ is
uniquely represented, in the original frame, by
\begin{equation}\label{eq:section}
 s=w^q\exp((q/r)A_K(w))f(w),\qquad
 |s|_{h^m}^2=|f(w)|^2e^{-(q/r)\Psi_K(w)},
\end{equation}
where $f$ is holomorphic at zero and has only nonnegative Taylor powers.
The volume in this chart is $R^2(1+O(R))\dA_w$. Put
$\eta=\sqrt{rq}\,v$. On each fixed $|\eta|\leq L$,
\[
 (q/r)\Psi_K=2\eta^2+o(1).
\]
For an observation $\eta_0$ in a fixed bounded interval, restriction to
this annulus bounds the kernel from above. Its limiting diagonal
monomial norms, for $f(w)=\sum_{j\geq0}b_jw^j$, are
\[
 n_{j,L}=\frac{2\pi R^2}{\sqrt{rq}}I_L(j/\sqrt{rq}),\qquad
 I_L(s)=\int_{-L}^L e^{2s\eta-2\eta^2}\,\mathrm d\eta.
\]
After normalization by $rq/(\pi R^2)$, the limiting upper kernel is
\begin{equation}\label{eq:hardy}
 \frac12e^{-2\eta_0^2}\int_0^\infty\frac{e^{2s\eta_0}}{I_L(s)}\,\mathrm ds.
\end{equation}
The sum-to-integral passage uses mesh $1/\sqrt{rq}$. To control its tail,
integrate $I_L$ on a fixed exterior subinterval beyond every observation
$\eta_0$; the resulting exponential bound is uniform in the mesh.
Thus Riemann sums and then $L\to\infty$ are justified. Since
$I_\infty(s)=\sqrt{\pi/2}\,e^{s^2/2}$, \eqref{eq:hardy} tends to
\begin{equation}\label{eq:gaussian}
 \frac1{\sqrt{2\pi}}\int_0^\infty e^{-(s-2\eta_0)^2/2}\,\mathrm ds
 =\PhiG(2\eta_0).
\end{equation}

For the lower bound fix $T$ and use \eqref{eq:section} with
$f$ in the span of $w^j$, $0\leq j\leq T\sqrt{rq}$. Cut off at
$1+\delta$, taking the cutoff equal to one through $1+\delta/2$.
The cutoff annulus lies outside the true hole. Equations
\eqref{eq:ringK} and \eqref{eq:section} bound the squared source norm,
as a quadratic form on the entire span, by
\[
 C(R\delta)^{-2}e^{-cq\delta^2}
 \quad\hbox{times the trial Gaussian norm}.
\]
The linear term $2jv$ is absorbed because $\sqrt q\,\delta\to\infty$.
Lemma~\ref{lem:dbar} gives a squared correction norm at most
\begin{equation}\label{eq:smallcorr}
 C(q\delta^2)^{-1}e^{-cq\delta^2}
 \quad\hbox{times that norm}.
\end{equation}
These estimates contain no polynomial factor in $m$.

The trial Gram form converges in relative operator norm to the diagonal
form
\begin{equation}\label{eq:gram}
 n_j=\frac{2\pi R^2}{\sqrt{rq}}\sqrt{\frac\pi2}
          \exp(j^2/(2rq)).
\end{equation}
Indeed the weight ratio is uniform on each fixed $\eta$-band. Outside
it, radial domination and angular Parseval bound arbitrary combinations
by diagonal tails $e^{-c\eta^2+2T|\eta|}$. First pass to the limit on
the fixed band and then enlarge it. Interior compact gaps and the
logarithmic singularity provide the remaining domination. This controls
the whole quadratic form without a dimension factor.

At points of radius comparable to $R$, a local submean estimate gives
\begin{equation}\label{eq:eval}
 |v(z)|_{h^m}^2\leq C(q/R^2)\|v\|_{h^m}^2
\end{equation}
for a section holomorphic near the point on a disk of radius
$R/\sqrt q$. Remove the harmonic linear and quadratic Taylor terms
there. The scaled remaining quadratic coefficient is bounded and the
third-order error is $O(q^{-1/2})$, which proves the estimate.
It applies to the corrections, since their sources lie on the distant
cutoff annulus. Together with \eqref{eq:smallcorr}, it shows that the
corrected global trial kernel differs by $o(q/R^2)$ from the local trial
kernel. The normalized lower limit is \eqref{eq:gaussian} with upper
limit $T$ instead of infinity; let $T\to\infty$.

At a nearest boundary point,
\[
 v=\frac{d_t(z)}{R|F_R'(e^{i\theta})|}
       +O(d_t(z)^2/R^2)+o(q^{-1/2}),\qquad
 m\rho(p)R^2|F_R'(e^{i\theta})|^2=rq(1+O(R)).
\]
Also $m\rho(z)/a(z)=(rq/R^2)(1+o(1))$ in the fixed transition band.
These identities convert $2\eta_0$ to the argument in \eqref{eq:edge}
and give its exact density normalization.

\section{The full kernel and saturation}
At points $|z|\asymp R$ the full, unselected kernel obeys
\begin{equation}\label{eq:full}
 B_m^{\rm full}(z)=\frac{m\rho(z)}{\pi a(z)}(1+o(1)),
\end{equation}
uniformly in an $o(R)$ neighborhood of $\Gamma_t$. Here is a direct
localization argument, which also supplies the saturation part of
Theorem~\ref{thm:corrected}. Remove the harmonic Taylor terms of $m\phi$
at the observation and rescale by $(m\rho(z))^{-1/2}$. On every fixed
rescaled disk the remaining weight tends to $|u|^2$ and the volume to
the constant $a(z)$. Taylor's estimate is uniform because the third
derivatives at radius $R$ are $O(R^{2r-3})$, so the normalized third
order error is $O(q^{-1/2})$. Restriction to such disks bounds the
evaluation norm from above; their Gaussian constant-function norm tends
to $\pi$ as their radius grows. This proves the upper bound in
\eqref{eq:full}.

For the lower bound use the holomorphic exponential cancelling the
harmonic Taylor terms, and cut it off on a disk of radius $\varepsilon R$
centered at the observation, where $\varepsilon\to0$ and
$\sqrt q\,\varepsilon\to\infty$. Its local squared norm is
$(\pi a(z)/(m\rho(z)))(1+o(1))$ relative to its unit weighted value
at the center. Quadratic lower domination bounds the source annulus by
an exponential $e^{-cq\varepsilon^2}$. Add the same bounded bump of
scale $R$ to the original metric and use the pointwise $\bar\partial$
estimate. The relative squared correction is bounded by
$C(q\varepsilon^2)^{-1}e^{-cq\varepsilon^2}=o(1)$, after decreasing
$c$ if needed. Its value at the center is $o(1)$ by the local submean
estimate \eqref{eq:eval}. This proves the matching lower bound.

For interior saturation, use the polynomial chart above and put
$u=1-|w|>0$. Restrict the norm in \eqref{eq:section} to the exterior
half-band $0\leq v\leq c/\sqrt q$. Every nonnegative Taylor monomial
of $f$ has norm at least $cR^2/\sqrt q$. Angular Parseval and
Cauchy--Schwarz give
\[
 B_{m,q}(z)\leq
 C\frac{\sqrt q}{R^2}\frac{e^{-cqu^2}}{1-|w|^2}.
\]
If $u\sqrt q\to\infty$, division by $q/R^2$ makes this tend to zero.
The chart error is $O(q^{-10})$ and does not alter the hypothesis or
the conclusion.

For exterior saturation, use a Gaussian peak at $z$ and cut it off on
a disk of radius
\[
 \varepsilon R,\qquad
 \varepsilon=\min\{d_t(z)/(3R),q^{-1/4}\}.
\]
It is wholly outside the hole. Again $\varepsilon\to0$ and
$\sqrt q\varepsilon\to\infty$. Apply Lemma~\ref{lem:dbar} using the
singular obstacle metric, which equals the original metric on this
source. The correction has negligible norm and value and vanishes to
order $q$ at $x$. Thus a member of $V_{m,q}$ attains the full Gaussian
density asymptotically. Equation~\eqref{eq:full} is its upper bound.
This proves saturation to one and finishes Theorem~\ref{thm:corrected}.

\section{The universal range and compact sharpness}
By Lemma~\ref{lem:hole}, $c_R-1=O(R)$ in every fixed norm required
here. The observation in \eqref{eq:raw} satisfies
\[
 |z_m|/R=1+\tau/(2\sqrt{rq})+O(q^{-1}).
\]
Since the boundary slope tends to zero, its corrected normal argument
is $\tau+o(1)$ whenever $\sqrt q R\to0$. This condition is precisely
$q=o(m^{1/(r+1)})$. Moreover $a(z)=1+o(1)$ and
$\rho(z)=r^2|z|^{2r-2}(1+o(1))$ there, so the two normalizations agree.
The edge theorem proves universal sufficiency.

The same reasoning proves the sufficiency of \eqref{eq:criterion}.
For necessity choose, on a subsequence, an angle realizing the maximum
boundary displacement and observe at $|z|=R$ with $\tau=0$. A nonzero
finite subsequential scaled displacement gives a nonzero Gaussian
argument; an infinite one gives saturation to zero or one. Radial and
normal distances have relative error tending to zero, since the graph
has vanishing slope. In either case the limit differs from $1/2$.

We give a compact model for sharpness without assumptions at infinity.
On $\CP^1$ choose a radial metric on $\mathcal O(d)$ with local potential
$f(u)$, $u=|w|^2$, equal to $u^r$ near zero, and with
$g(u)=uf'(u)$ strictly increasing on $(0,\infty)$ and tending to $d$.
For large $u$ require $g(u)=d-\epsilon/u$, with $\epsilon>0$.
Such a smooth interpolation exists for a sufficiently large integer
$d$: prescribe a positive $g'$ equal to $r^2u^{r-1}$ near zero and
$\epsilon/u^2$ near infinity, with total integral $d$.
Then $f=d\log u+\mathrm{constant}+\epsilon/u$ near infinity; the
infinity-frame potential is smooth with strictly positive curvature.
Choose a smooth radial volume equal to $\dA_w$ near zero.

For sufficiently small $t$, the radial hole lies in the exact-model
region and is $|w|<R$: rotational invariance and uniqueness make it a
disk, and its curvature mass is $rR^{2r}=t$. The corrected transition
and saturation theorem therefore apply to this compact example.

Change the allowed observation coordinate to $w=z+c z^2$, $c>0$.
The same compact metric now has
$\phi(z)=|z|^{2r}+O(|z|^{2r+1})$, and the volume still has $a(0)=1$.
Vanishing order is invariant. At the raw observation $z=R$,
\[
 2\sqrt{mr^2R^{2r-2}}\,(|w|-R)
 =2c\sqrt r\,\sqrt q R.
\]
The two curvature-model denominators have ratio tending to one.
Every excluded sequence has a subsequence on which $\sqrt q R$ tends
to a positive finite number or to infinity. Theorem~\ref{thm:corrected}, using its saturation conclusion when the
argument diverges, then gives a value strictly larger than $1/2$. This proves sharpness
and completes Theorem~\ref{thm:range}.

\section{Corrections from the smooth weight}\label{sec:corrections}
Write
\[
 \phi(z)=|z|^{2r}+P_{2r+1}(z,\bar z)+O(|z|^{2r+2}),\qquad
 P(\theta)=P_{2r+1}(e^{i\theta},e^{-i\theta}),
\]
where the polynomial is real homogeneous. Let $\Lambda$ be the
circle Dirichlet-to-Neumann operator,
$\Lambda e^{in\theta}=|n|e^{in\theta}$. The first boundary coefficient is
\begin{equation}\label{eq:b1}
 c_R=1+Rb_1+O(R^2),\qquad
 b_1=-\frac{((2r+1)-\Lambda)P}{4r^2}.
\end{equation}
To derive it, write the interior obstacle to first order as
$1+2r\log s+R H$, with harmonic $H$. Value matching on
$s=1+Rb_1$ gives $H|_{s=1}=P$; the terms involving $b_1$ cancel since
the radial reference already matches first derivatives. Derivative
matching gives $4r^2b_1+(2r+1)P-\Lambda P=0$.
Pure harmonic terms disappear, as required by invariance under frame
changes.

For $q=O(m^{1/(r+1)})$, at the raw observation the corrected argument is
\begin{equation}\label{eq:firstshift}
 \tau+\frac{\sqrt q R}{2r^{3/2}}((2r+1)-\Lambda)P(\theta)+o(1),
\end{equation}
uniformly in $\theta$. The kernel differs from $\PhiG$ of the displayed
argument by $o(1)$. The error is small because $\sqrt q R^2\to0$ in
this range. For $w=z+cz^2$, $P=2r\Rea(ce^{i\theta})$ and the shift is
$2\sqrt{rq}R\Rea(ce^{i\theta})$, agreeing with the compact example.

For every fixed $J$ the parameter construction gives
\[
 c_R=1+\sum_{j=1}^J R^j b_j+O(R^{J+1})
\]
in every fixed norm needed. There is an explicit recursion: expand the
interior harmonic obstacle in powers of $R$. At order $j$, value
matching determines the boundary data of its new harmonic term from
the degree $2r+j$ Taylor term and all preceding coefficients. Harmonic
extension determines this term. Derivative matching is then
$4r^2b_j+\text{known forcing}=0$. The new $b_j$ does not enter the
order-$j$ value equation. Nonlinear lower-order terms contribute when
$j\geq2$.

The truncated curve changes the physical center by $O(R^{J+2})$.
Using it in the geometric correction is valid whenever
\[
 \sqrt q R^{J+1}\to0,
 \quad\hbox{equivalently}\quad
 q=o(m^{(J+1)/(r+J+1)}).
\]
If $b_1=\cdots=b_{j-1}=0$ but $b_j\not\equiv0$, the exact criterion
\eqref{eq:criterion} makes the fixed-weight uncorrected range
$q=o(m^{j/(r+j)})$, sharply in the all-direction sense. Higher terms
are relatively $O(R)$ at an angle where $b_j\ne0$.

In the radial case $\phi=f(u)$, $u=|z|^2$, the hole is a disk by
uniqueness and rotational invariance. Its exact center satisfies
\[
 m u_*f'(u_*)=q.
\]
The asymptotic mode variance is $m u\rho(u)$, with
$\rho(u)=f'(u)+uf''(u)$. The corresponding Gaussian argument is
\begin{equation}\label{eq:radialarg}
 \frac{m u f'(u)-q}{\sqrt{m u\rho(u)}}.
\end{equation}
For bounded arguments, Taylor expansion at $u_*$ identifies it with
$2\sqrt{m\rho(u_*)}(\sqrt u-\sqrt{u_*})+o(1)$, so
Theorem~\ref{thm:corrected} proves this radial formula for all sublinear
orders. For $f=u^r+c u^{r+1}+\cdots$, its first raw-center shift is
$c(r+1)r^{-3/2}\sqrt q R^2$, with critical growth $m^{2/(r+2)}$.
This larger radial threshold does not describe arbitrary coordinates.

No fixed finite Taylor jet determines the correction for every
$q=o(m)$. For example, near zero take
\[
 f(u)=u^r+c u^r e^{-1/\sqrt u},\qquad c>0,
\]
smoothly extended at zero and interpolated globally as above.
It has the same Taylor jet as $u^r$ at every finite order, but at
$u=R^2$ its relative contribution to $uf'(u)$ is
$c e^{-1/R}(1+1/(2rR))$. For $q=\lfloor m/\log m\rfloor$,
$\sqrt q e^{-1/R}/R\to\infty$. The raw observation with $\tau=0$
is outside the true center by many transition widths, giving limit one.
The exact quadrature-domain correction includes this flat contribution.

\section{Scope and checks}
This paper resolves the original compact-curve problem P22, including
uniformity in direction and every fixed integer $r\geq1$. Its general
correction uses the actual smooth curvature, rather than only polynomial
examples. It does not solve higher-dimensional coupled anisotropic
problems or arbitrary fiber-selected singularities. Kernel positivity,
the order-preserving singular correction and the explicit compact
models account for eventual nonzero selected spaces.

The manuscript's previously established smaller range
$q=o(m^{1/(2r+1)})$ comes from requiring $mR^{2r+1}=o(1)$ on a disk.
The argument removes the boundary harmonic error and uses its
normal variation over a width $R/\sqrt q$, leading to $\sqrt q R$.
No priority claim or exhaustive literature claim is made.

\begin{thebibliography}{99}
\bibitem{catalogue}
\emph{Open problems in non-K\"ahler pluripotential theory: full statements,
known results, and remaining gaps}, catalogue accompanying Mingchen Xia's
monograph, 2 October 2026, supplied as \texttt{document.pdf}, P22.
\bibitem{book}
M. Xia, \emph{Singularities in Hermitian pluripotential theory}, public
working manuscript, Theorem 26.13.1 and the growing-jet question.
\url{https://mingchenxia.github.io/Papers/SHPT_vibe.pdf}.
\bibitem{D} J.-P. Demailly,
\emph{Multiplier ideal sheaves and analytic methods in algebraic geometry},
Theorem (5.2) and the singular-metric extension following it.
\url{https://www-fourier.univ-grenoble-alpes.fr/~demailly/manuscripts/trieste.pdf}.
\bibitem{DA} J.-P. Demailly,
\emph{Complex analytic and differential geometry}, Chapter VI, Section 10.
\url{https://www-fourier.univ-grenoble-alpes.fr/~demailly/manuscripts/agbook.pdf}.
\bibitem{K} S. Kumaresan,
\emph{Implicit function theorem for Banach spaces}, Theorem 2.
\url{https://4dspace.mtts.org.in/expository-article-download.php?ai=78}.
\bibitem{WD} Y. Wang and J. Du,
\emph{Boundedness of Cauchy singular integral operator under H\"older norm},
Zeitschrift f\"ur Analysis und ihre Anwendungen \textbf{43} (2024),
209--235, Theorem 1.2.
\url{https://ems.press/content/serial-article-files/47437}.
\bibitem{W} J. Wehler,
\emph{Riemann surfaces}, lecture notes, 30 May 2019, Corollary 13.12.
\url{https://www.mathematik.uni-muenchen.de/~wehler/20190530_RiemannSurfacesScript.pdf}.
\bibitem{F} F. Forstneri\v c,
\emph{Analytic disks with boundaries in a maximal real submanifold of
$\mathbb C^2$}, Theorem 1 and Remark 1.
\url{https://numdam.org/item/10.5802/aif.1076.pdf}.
\bibitem{L} C. LeBrun,
\emph{Twistors, holomorphic disks, and Riemann surfaces with boundary}.
\url{https://arxiv.org/abs/math/0508038}.
\bibitem{RWN} J. Ross and D. Witt Nystr\"om,
\emph{The Hele-Shaw flow and moduli of holomorphic discs}.
\url{https://arxiv.org/abs/1212.2337}.
\bibitem{SS} S. Serfaty and J. Serra,
\emph{Quantitative stability of the free boundary in the obstacle problem},
Analysis \& PDE \textbf{11} (2018), 1803--1849.
\url{https://doi.org/10.2140/apde.2018.11.1803}.
\end{thebibliography}
\end{document}
